\section{回帰検査}
\label{sec:checks}

この章は、\texttt{Regression.lean} が例として証明している数値範囲を、シミュレーターの単体テストへ写す。範囲そのものが新しい物理の主張ではない。テストは、主張の種類ごとにファイルを分ける。内在角のテストは \texttt{written\_at} だけを呼ぶ。層のテストは力の関数だけを呼ぶ。チャート比較のテストは、そのどちらも呼ばない。

\subsection{書かれた流れ}

次を機械的に検査する。

\begin{itemize}
\item \(\delta(t)=\delta_0+\nu t\) かつ \(\ddot\delta=0\)。
\item \(\sigma(t)\) の三次式と、\(\ddot\sigma=-2m\delta(t)\)。
\item \(\phi=(\sigma+\delta)/2\)、\(\theta=(\sigma-\delta)/2\)。
\item 数点の時刻で \(E(t)=E(0)=\frac12\nu\dot\sigma_0+\frac m2\delta_0^2\)。
\item \(p_\sigma(t)=\nu/2\) が一定。
\item サンプル \((m,\phi,\theta,\ddot\phi,\ddot\theta)=(1,1,0,-1,-1)\) は \texttt{PaperActualEL} を満たし、\(\ddot\delta=-m\delta\) を満たさない。
\item 同じサンプルの近傍 \((1,1,0,1,-1)\) は \texttt{PaperClaimedEL} を満たし、\texttt{PaperActualEL} を満たさない。
\item \(m\ge 0\) に対し \(\sqrt{2m}=m\) となるのは \(m=0,2\) だけである。この検査は振動子モジュールに置き、書かれた流れのテストから参照しない。
\end{itemize}

\subsection{層と力}

\begin{itemize}
\item \(\gamma_s^{\mathrm{eq}}(0)=1\)。
\item 導関数が \(-2\cosh x\sin x\) と、選んだ差分幅の中で一致する。
\item 見つけた第一正根が \(\pi/4\) と \(1\) のあいだにあり、その点で \(\gamma_s^{\mathrm{eq}}=0\)。
\item \(\kappa>0\)、\(r>0\)、\(\gamma<0\) で \texttt{gravityLayerTerm}\(>0\)。\(\gamma>0\) で項は負。
\item 二体の \texttt{gravityLayerPush} が反対称。
\item 斥力層で、外向き単位分離との内積が正。
\item \(N=3\) の和が、自分自身を項に含まない。
\item \(\gamma_G,\gamma_E,r\neq 0\) のとき、\texttt{pairChannelForce}\(=0\) と \(\kappa\gamma_E=\mathrm{coul}\,\gamma_G\) が同値。サンプルは半径を二種類に取り、零の有無が半径で変わらないことを見る。
\end{itemize}

\subsection{チャート比較}

\texttt{Regression} は次の区間を例として持つ。

\begin{itemize}
\item 水星の近日点余剰について \(42<\texttt{arcsecPerCentury}<44\)。
\item 太陽による光線の偏折について \(17/10<\texttt{deflectionArcsec}<9/5\)。
\end{itemize}

これらの数は、外部の軌道要素と SI スタンドインから出るチャート量である。内在角の閉形式の出力と比べない。テスト名に \texttt{chart} を入れる。

\subsection{核の長さが入る範囲}

散乱の符号が反転する読みでは、第一ノードの半径 \(r_1\) が
\[
\frac12<r_1<1
\]
（fm の入力として渡す）にあるとき、対応する長さ
\[
\texttt{lengthOfReversal}(r_1)=2\cdot\texttt{resonanceRoot1}\cdot r_1
\]
が
\[
\frac{14}{15}<\texttt{lengthOfReversal}(r_1)<\frac{19}{10}
\]
を満たす。これは、渡した \(r_1\) に対する長さの検査である。力を時間積分した出力ではない。\(\alpha_s\) をこの範囲から逆算しない。

強い力の結合定数を返す関数、核融合率を返す関数は、このリポジトリの記録に無い。テストもそれらを期待値にしない。
